\documentclass[11pt]{article}
\usepackage[T1]{fontenc}
\usepackage{lmodern}
\usepackage[sc]{mathpazo}
\usepackage{amsmath,amssymb,amsthm,mathtools}
\usepackage{microtype}
\usepackage[letterpaper,top=0.85in,bottom=0.87in,left=0.95in,right=0.95in,headheight=15pt,headsep=19pt]{geometry}
\usepackage{booktabs,array,tabularx,enumitem}
\usepackage{xcolor,tikz,titlesec,fancyhdr,listings}
\usetikzlibrary{arrows.meta,calc,positioning,decorations.pathreplacing}
\definecolor{ink}{HTML}{173B4D}
\definecolor{teal}{HTML}{277D79}
\definecolor{gold}{HTML}{B58B4A}
\definecolor{slate}{HTML}{5A6971}
\definecolor{mist}{HTML}{EEF4F4}
\usepackage[colorlinks=true,linkcolor=ink,citecolor=teal,urlcolor=teal,breaklinks=true]{hyperref}
\usepackage[nameinlink,noabbrev]{cleveref}
\crefname{theorem}{Theorem}{Theorems}
\crefname{lemma}{Lemma}{Lemmas}
\crefname{proposition}{Proposition}{Propositions}
\crefname{corollary}{Corollary}{Corollaries}
\crefname{definition}{Definition}{Definitions}
\crefname{section}{Section}{Sections}
\crefname{figure}{Figure}{Figures}
\crefname{table}{Table}{Tables}
\hypersetup{pdfauthor={John Tripodo},pdftitle={Separated distances in fixed dimension},pdfsubject={Two-circle tolerance family and a dimension-29 counterexample}}
\numberwithin{equation}{section}
\newtheorem{theorem}{Theorem}[section]
\newtheorem{proposition}[theorem]{Proposition}
\newtheorem{lemma}[theorem]{Lemma}
\newtheorem{corollary}[theorem]{Corollary}
\theoremstyle{definition}\newtheorem{definition}[theorem]{Definition}
\theoremstyle{remark}\newtheorem{remark}[theorem]{Remark}
\newcommand{\R}{\mathbb R}\newcommand{\N}{\mathbb N}
\newcommand{\ip}[2]{\langle #1,#2\rangle}
\newcommand{\norm}[1]{\lVert #1\rVert}
\newcommand{\opnorm}[1]{\lVert #1\rVert_{\mathrm{op}}}
\newcommand{\T}{\mathsf T}\newcommand{\one}{\mathbf1}
\newcommand{\code}[1]{{\ttfamily\small #1}}
\newcommand{\TripodoCommit}{\texttt{1e73a12}}
\DeclareMathOperator{\diam}{diam}\DeclareMathOperator{\diag}{diag}
\titleformat{\section}{\Large\bfseries\color{ink}}{\thesection}{.7em}{}
\titleformat{\subsection}{\normalsize\bfseries\color{ink}}{\thesubsection}{.65em}{}
\titlespacing*{\section}{0pt}{22pt}{10pt}
\titlespacing*{\subsection}{0pt}{16pt}{6pt}
\setlength{\parindent}{1.15em}
\setlength{\parskip}{2pt}
\setlength{\tabcolsep}{5pt}
\renewcommand{\arraystretch}{1.12}
\linespread{1.055}
\emergencystretch=2em
\setlength{\abovecaptionskip}{7pt}
\setlength{\belowcaptionskip}{5pt}
\pagestyle{fancy}\fancyhf{}
\fancyhead[L]{\scriptsize\scshape\color{slate}Separated distances in fixed dimension}
\fancyhead[R]{\scriptsize\color{slate}John Tripodo}
\fancyfoot[L]{\scriptsize\color{slate}20 October 2026}
\fancyfoot[R]{\small\color{ink}\thepage}
\renewcommand{\headrulewidth}{0.3pt}
\renewcommand{\footrulewidth}{0pt}
\lstset{basicstyle=\ttfamily\scriptsize,keywordstyle=\color{ink}\bfseries,commentstyle=\color{slate},stringstyle=\color{teal},breaklines=true,columns=fullflexible,keepspaces=true,showstringspaces=false,frame=single,rulecolor=\color{teal!35},backgroundcolor=\color{mist},xleftmargin=5pt,xrightmargin=5pt,aboveskip=10pt,belowskip=10pt}
\begin{document}
\thispagestyle{fancy}
\noindent{\sffamily\small\bfseries\color{teal}METRIC GEOMETRY\quad /\quad ERD\H{O}S PROBLEM 670}\par
\vspace{18pt}
\noindent{\fontsize{30}{34}\selectfont\bfseries\color{ink}Separated distances\\[3pt]in fixed dimension}\par
\vspace{10pt}
\noindent{\Large\color{slate}Two circles and a\\[2pt]dimension-29 counterexample}\par
\vspace{15pt}
\noindent{\large John Tripodo}\hfill{\small 20 October 2026}\par
\vspace{3pt}
\noindent{\small\color{slate}\href{https://orcid.org/0009-0006-9180-4953}{ORCID 0009-0006-9180-4953}}\par
\vspace{13pt}\noindent{\color{gold}\rule{\textwidth}{1pt}}\par
\vspace{8pt}
\noindent\textbf{Abstract.}
A finite Euclidean point set is distance-separated when the lengths of any
two different unordered edges differ by at least one. The counting bound
for its diameter is asymptotically $n^2/2$, whereas Erd\H{o}s asked for a
coefficient-one lower bound. We give a two-circle construction: for every
$0<\varepsilon\le1/2$ and every $n\ge4\lceil8192/\varepsilon\rceil$, there
is a distance-separated $n$-point set in
$\R^{10\lceil8192/\varepsilon\rceil-1}$ with diameter
$<(1/2+\varepsilon)n^2$. A uniform form has coefficient
$1/2+102400/d+40960/n$. Independently, an explicit six-vertex macro seed
gives arbitrarily large configurations in the single space $\R^{29}$
with diameter $<0.998n^2$, refuting the maintained fixed-dimensional
coefficient-one statement. A short-prime-interval theorem also gives
diameter $<0.999n^2$ in $\R^{29}$ for every $n\ge400000$. Two orthogonal circle channels separate local
ruler gaps from cross-edge digits; the latter have an exact monotone
square-root formula. The small-dimensional proof uses three rational
matrix certificates valid for all admissible finite cluster sizes.
The two-circle family and the prime-first dimension-29 theorem have
corresponding Lean~4 proofs; the all-size corollary is a written consequence. The general family
uses uniform colex estimates; the dimension-29 construction replaces them
with a finite rational certificate.
\vspace{12pt}
\begin{center}
\fcolorbox{teal!30}{mist}{\parbox{.92\textwidth}{\small\color{ink}
\textbf{GitHub manuscript, 20 October 2026.}
The prime-first theorem and its all-size corollary have different bounds
and verification scopes. The near-counting family fixes the dimension
separately for each tolerance.}}
\end{center}
\vspace{5pt}
\noindent{\footnotesize\textbf{Keywords:} separated distances, Euclidean diameter,
Gram completion, quadratic-residue rulers, Lean formalization.}
\clearpage
